% =============================================================================
% PNDR Survey - Systematic Literature Review
% =============================================================================
% Article: Avaliação dos Instrumentos da PNDR: Uma Revisão Sistemática
%          da Literatura Empírica (2005--2025)
%
% Author:  Hermelino Nepomuceno de Souza
% =============================================================================
\documentclass[12pt,a4paper]{article}

% --- Encoding and Language ---
\usepackage[utf8]{inputenc}
\usepackage[T1]{fontenc}
\usepackage[brazil]{babel}

% --- Page Layout ---
\usepackage[top=2.5cm, bottom=2.5cm, left=3cm, right=2.5cm]{geometry}
\usepackage{setspace}
\onehalfspacing

% --- Fonts ---
\usepackage{lmodern}

% --- Math ---
\usepackage{amsmath,amssymb}

% --- Tables ---
\usepackage{booktabs}
\usepackage{tabularx}
\usepackage{longtable}
\usepackage{multirow}

% --- Figures ---
\usepackage{graphicx}
\graphicspath{{../figures/}}
\usepackage{float}

% --- Citations ---
\usepackage{natbib}
\bibliographystyle{apalike}

% --- Hyperlinks ---
\usepackage{hyperref}
\hypersetup{
    colorlinks=true,
    linkcolor=blue,
    citecolor=blue,
    urlcolor=blue
}

% --- Misc ---
\usepackage{enumerate}
\usepackage{caption}
\usepackage{subcaption}

% =============================================================================
\title{Avaliação dos Instrumentos da PNDR: Uma Revisão Sistemática da Literatura Empírica (2005--2025)}

\author{
    Hermelino Nepomuceno de Souza\thanks{CAEN/UFC. E-mail: hermelino@example.com}
}

\date{\today}

% =============================================================================
\begin{document}

\maketitle

% --- Abstract ---
\begin{abstract}
    % TODO: Inserir resumo
    \noindent\textbf{Palavras-chave:} PNDR; Revisão Sistemática; Fundos Constitucionais; Fundos de Desenvolvimento; Incentivos Fiscais.

    \vspace{0.5cm}
    \noindent\textbf{JEL:} R11; R58; H50; O18.
\end{abstract}

\newpage

% =============================================================================
% SECTIONS
% =============================================================================

\section{Introdução}
\label{sec:introducao}
% TODO: Migrar conteúdo de 1-1-introducao.tex

\section{Política Regional no Brasil}
\label{sec:politica-regional}
% TODO: Migrar conteúdo de 1-2-politica-regional-no-brasil.tex

\section{Método}
\label{sec:metodo}
% TODO: Migrar conteúdo de 1-3-metodo.tex

\section{Discussão dos Resultados}
\label{sec:resultados}
% TODO: Migrar conteúdo de 1-4-discussao-dos-resultados.tex

\section{Considerações Finais}
\label{sec:consideracoes}
% TODO: Migrar conteúdo de 1-5-consideracoes-finais.tex

% =============================================================================
% REFERENCES
% =============================================================================
% \bibliography{references}

\end{document}
